% ========================================
% PACKAGES DE BASE
% ========================================
\usepackage[utf8]{inputenc}                    % Encodage UTF-8 pour les caractères accentués
\usepackage[T1]{fontenc}                       % Encodage de sortie pour les polices européennes
\usepackage[french]{babel}                     % Support de la langue française
\usepackage{microtype}
% ========================================
% PACKAGES MATHÉMATIQUES
% ========================================
\usepackage{mathtools}                         % Extension d'amsmath avec outils supplémentaires
\usepackage{amsfonts}                          % Polices mathématiques AMS
\usepackage{amssymb}                           % Symboles mathématiques AMS
\usepackage{amsthm}                            % Environnements de théorèmes
\usepackage{mathrsfs}                          % Lettres calligraphiques \mathscr{}

% ========================================
% PACKAGES GRAPHIQUES ET COULEURS
% ========================================
\usepackage{graphicx}                          % Inclusion et manipulation d'images
\usepackage[dvipsnames]{xcolor}                % Couleurs étendues (noms prédéfinis)
\usepackage{tikz}                              % Graphiques TikZ
\usepackage{tikz-cd}
\usetikzlibrary{automata,positioning,arrows.meta}                           % Diagrammes commutatifs
\usetikzlibrary{calc}

% ========================================
% PACKAGES DE MISE EN PAGE
% ========================================
\usepackage[a4paper, margin=1in]{geometry}    % Configuration des marges
\usepackage{fancyhdr}                          % En-têtes et pieds de page personnalisés
\usepackage{titlesec}                          % Personnalisation des titres de sections
\usepackage{setspace}                          % Contrôle de l'interligne
\usepackage{titling} %Récupération du nom des auteurs                            
% Configuration détaillée des marges
\geometry{
    top=2.5cm,
    bottom=2.5cm,
    left=2.5cm,
    right=2.5cm,
    headheight=15.03pt,
    headsep=20pt,
    footskip=30pt
}

% ========================================
% PACKAGES DE MISE EN FORME
% ========================================
\usepackage{textcomp}                          % Symboles textuels supplémentaires
\usepackage{soul}                              % Soulignement et surlignage avancés
\usepackage{mdframed}                          % Boîtes encadrées personnalisables
\usepackage{tcolorbox}                         % Boîtes colorées avancées

% ========================================
% LISTES ET ÉNUMÉRATIONS
% ========================================
\usepackage[shortlabels]{enumitem}
\setlistdepth{9}
\renewlist{enumerate}{enumerate}{9}
\renewlist{itemize}{itemize}{9}

% Configuration des niveaux d'énumération
\setlist[enumerate,1]{label=(\roman*)}
\setlist[enumerate,2]{label=(\alph*)}
\setlist[enumerate,3]{label=\arabic*.}
\setlist[enumerate,4]{label=(\roman*)}
\setlist[enumerate,5]{label=(\alph*)}

% Configuration des niveaux d'itemize
\setlist[itemize,1]{label=$\bullet$}
\setlist[itemize,2]{label=$\circ$}
\setlist[itemize,3]{label=$-$}
\setlist[itemize,4]{label=$\ast$}

% ========================================
% HYPERLIENS
% ========================================
\definecolor{linkcolor}{RGB}{0,0,139}
\definecolor{citecolor}{RGB}{139,0,0}
\definecolor{urlcolor}{RGB}{0,100,0}

\usepackage[
    colorlinks=true,
    linkcolor=linkcolor,
    citecolor=citecolor,
    urlcolor=urlcolor,
    bookmarks=true,
    bookmarksnumbered=true,
    pdfstartview=FitH
]{hyperref}

% ========================================
% POLICES MATHÉMATIQUES
% ========================================
\usepackage{fouriernc}                         % Police Fourier
\DeclareMathAlphabet{\mathbb}{U}{msb}{m}{n}    % \mathbb standard (AMS)
\DeclareMathAlphabet{\mathcal}{OMS}{cmsy}{m}{n} % \mathcal standard

% ========================================
% STYLES POUR LES BOÎTES
% ========================================
\definecolor{lime}{RGB}{0,0,128}

% Macro pour définir des styles rapidement
\newcommand{\defstyle}[2]{\mdfdefinestyle{#1}{
    leftmargin=0pt,
    rightmargin=0pt,
    innerleftmargin=10pt,
    innerrightmargin=10pt,
    innertopmargin=0pt,
    innerbottommargin=2pt,
    leftline=true,
    rightline=false,
    topline=false,
    bottomline=false,
    linecolor=#2,
    linewidth=2pt,
    backgroundcolor=white
}}

% Styles spécifiques
\defstyle{definitionstyle}{RoyalBlue}
\defstyle{theoremstyle}{lime}
\defstyle{attentionstyle}{Red}

% Style classique (sans bordure colorée)
\mdfdefinestyle{classicstyle}{
    leftmargin=0pt,
    rightmargin=0pt,
    innerleftmargin=0pt,
    innerrightmargin=0pt,
    innertopmargin=0pt,
    innerbottommargin=2pt,
    leftline=true,
    rightline=false,
    topline=false,
    bottomline=false,
    linecolor=white,
    linewidth=2pt,
    backgroundcolor=white
}

% ========================================
% ENVIRONNEMENTS DE THÉORÈMES
% ========================================

% Définitions et axiomes
\newmdtheoremenv[style=definitionstyle,nobreak=true]{definition}{Définition}[section]
\newmdtheoremenv[style=definitionstyle,nobreak=true]{axiome}[definition]{Axiome}
\newmdtheoremenv[style=definitionstyle,nobreak=true]{conjecture}[definition]{Conjecture}
\newmdtheoremenv[style=definitionstyle,nobreak=true]{convention}[definition]{Convention}

% Théorèmes et résultats
\newmdtheoremenv[style=theoremstyle,nobreak=true]{theoreme}[definition]{Théorème}
\newmdtheoremenv[style=theoremstyle,nobreak=true]{lemme}[definition]{Lemme}
\newmdtheoremenv[style=theoremstyle,nobreak=true]{proposition}[definition]{Proposition}
\newmdtheoremenv[style=theoremstyle,nobreak=true]{corollaire}[definition]{Corollaire}

% Style amsthm
\newtheoremstyle{remarkstyle}
  {}{}
  {\normalfont}  % corps en droit
  {}
  {\itshape\bfseries}     % titre en italique
  {.}{ }{}

% Environnements auxiliaires
\theoremstyle{remarkstyle}


\newtheorem{remarque}[definition]{Remarque}
\surroundwithmdframed[style=classicstyle]{remarque}
% Style amsthm
\newtheoremstyle{normalstyle}
  {}{}
  {\normalfont}  % corps en droit
  {}
  {\normalfont\bfseries}     % titre en italique
  {.}{ }{}

% Environnements auxiliaires
\theoremstyle{normalstyle}
\newtheorem{exercice}[definition]{Exercice}
\surroundwithmdframed[style=classicstyle]{exercice}
\newtheorem{exemple}[definition]{Exemple}
\surroundwithmdframed[style=classicstyle]{exemple}
\newtheorem*{resolution}{Résolution}
\surroundwithmdframed[style=classicstyle]{resolution}
\newtheorem*{notation}{Notation}
\surroundwithmdframed[style=classicstyle]{notation}

% ========================================
% EN-TÊTES ET PIEDS DE PAGE
% ========================================
\pagestyle{fancy}
\fancyhf{}
\fancyhead[L]{\textit{\leftmark}}              % Chapitre à gauche
\fancyfoot[L]{\thepage}                        % Numéro de page à gauche
\fancyfoot[R]{\textit{\rightmark}}             % Section à droite

\renewcommand{\headrulewidth}{0pt}
\renewcommand{\footrulewidth}{0pt}

% Marques de chapitres et sections
\renewcommand{\chaptermark}[1]{%
    \markboth{\MakeUppercase{\chaptername\ \thechapter.\ #1}}{}}
\renewcommand{\sectionmark}[1]{%
    \markright{\MakeUppercase{\thesection.\ #1}}}

% Style pour les premières pages de chapitre
\fancypagestyle{plain}{%
    \fancyhf{}
    \fancyfoot[L]{\thepage}
    \fancyfoot[R]{\textit{\rightmark}}
}

% ========================================
% FORMAT DES TITRES DE SECTIONS
% ========================================
\titleformat{\chapter}[block]
    {\normalfont\LARGE\bfseries\scshape\color{RoyalBlue}\flushleft}
    {\thechapter.}
    {1em}
    {\scshape\color{RoyalBlue}}

\titleformat{\section}[block]
    {\normalfont\large\scshape\color{RoyalBlue}\flushleft}
    {\thesection}
    {1em}
    {}

\titleformat{\subsection}[block]
    {\normalfont\large\scshape\color{RoyalBlue}\flushleft}
    {\thesubsection}
    {1em}
    {}

\titleformat{\subsubsection}[block]
    {\normalfont\large\scshape\color{RoyalBlue}\flushleft}
    {\thesubsubsection}
    {1em}
    {}

% ========================================
% ENSEMBLES USUELS
% ========================================
\newcommand{\N}{\mathbb{N}}                    % Entiers naturels
\newcommand{\Z}{\mathbb{Z}}                    % Entiers relatifs
\newcommand{\Q}{\mathbb{Q}}                    % Rationnels
\newcommand{\R}{\mathbb{R}}                    % Réels
\newcommand{\C}{\mathbb{C}}                    % Complexes
\newcommand{\K}{\mathbb{K}}                    % Corps quelconque
\newcommand{\Sg}{\mathfrak{S}}                 % Groupe symétrique
\newcommand{\sn}{\mathbb{S}_n}                 % Sphère de dimension n

% ========================================
% LOGIQUE
% ========================================
\newcommand{\ssi}{\Leftrightarrow}             % Si et seulement si
\newcommand{\alors}{\Rightarrow}               % Implique
\newcommand{\si}{\Leftarrow}                   % Est impliqué par

% ========================================
% NORMES ET DISTANCES
% ========================================
\newcommand{\norm}[1]{\left\| #1 \right\|}     % Norme
\newcommand{\norme}[1]{\lVert #1 \rVert}       % Norme (variante)
\newcommand{\abs}[1]{\left| #1 \right|}        % Valeur absolue

% ========================================
% INTERVALLES
% ========================================
\newcommand{\interoo}[1]{\left] #1\right[}  % ]a,b[
\newcommand{\interof}[1]{\left] #1\right]}   % ]a,b]
\newcommand{\interfo}[1]{\left[ #1\right[}   % [a,b[
\newcommand{\interff}[1]{\left[ #1\right]}   % [a,b]

% ========================================
% TOPOLOGIE
% ========================================
\newcommand{\topologie}{\mathcal{T}}           % Topologie
\newcommand{\voisinage}{\mathcal{V}}           % Voisinages
\newcommand{\partie}{\mathscr{P}}              % Parties
\newcommand{\filtre}{\mathscr{F}}              % Filtre
\newcommand{\ouvert}{\mathcal{O}}              % Ouverts
\newcommand{\ferme}{\mathcal{F}}               % Fermés
\newcommand{\adh}{\operatorname{adh}}                % Adhérence
% ========================================
% FONCTIONS
% ========================================
\newcommand{\fonction}[3]{#1\colon\left\{\begin{array}{l}#2 \\#3\end{array}\right.}

% ========================================
% OPÉRATEURS MATHÉMATIQUES
% ========================================
\DeclareMathOperator{\dom}{dom}               % Domaine
\DeclareMathOperator{\im}{Im}                 % Image
\DeclareMathOperator{\e}{e}                   % Exponentielle
\DeclareMathOperator{\stab}{St}               % Stabilisateur
\DeclareMathOperator{\orb}{Orb}               % Orbite
\DeclareMathOperator{\coker}{coker}           % Conoyau
\DeclareMathOperator{\sgn}{sgn}               % Signature
\DeclareMathOperator{\Frac}{Frac}             % Corps des fractions
\DeclareMathOperator{\Eq}{Eq}                 % Égaliseur
\DeclareMathOperator{\Coeq}{Coeq}             % Coégaliseur
\DeclareMathOperator{\Cone}{Cone}             % Catégorie des cônes
\DeclareMathOperator{\colim}{colim}           % Colimite

% Limites avec indices
\newcommand{\limi}{\lim\limits}

% Intégrale différentielle
\newcommand{\intd}{\,\operatorname{d}}
\newcommand{\dpart}[2]{\frac{\partial#1}{\partial#2}}

% ========================================
% THÉORIE DES CATÉGORIES
% ========================================

% --- Catégories standards ---
\newcommand{\Set}{\mathbf{Set}}               % Ensembles
\newcommand{\Grp}{\mathbf{Grp}}               % Groupes
\newcommand{\Ab}{\mathbf{Ab}}                 % Groupes abéliens
\newcommand{\Ring}{\mathbf{Ring}}             % Anneaux
\newcommand{\CRing}{\mathbf{CRing}}           % Anneaux commutatifs
\newcommand{\Mod}{\mathbf{Mod}}               % Modules
\newcommand{\Vect}{\mathbf{Vect}}             % Espaces vectoriels
\newcommand{\Top}{\mathbf{Top}}               % Espaces topologiques
\newcommand{\Pos}{\mathbf{Pos}}               % Ensembles ordonnés
\newcommand{\Cat}{\mathbf{Cat}}               % Catégories
\newcommand{\Mon}{\mathbf{Mon}}               % Monoïdes

% --- Constructions ---
\newcommand{\Disc}{\mathbf{Disc}}             % Catégorie discrète
\newcommand{\op}{^{\mathrm{op}}}              % Catégorie opposée

% --- Objets et morphismes ---
\newcommand{\Ob}{\mathrm{Ob}}                 % Objets
\newcommand{\Hom}{\mathrm{Hom}}               % Morphismes
\newcommand{\End}{\mathrm{End}}               % Endomorphismes
\newcommand{\Aut}{\mathrm{Aut}}               % Automorphismes
\newcommand{\id}{\mathrm{id}}                 % Identité
\newcommand{\Id}{\mathrm{Id}}                 % Foncteur identité

% --- Groupes linéaires ---
\newcommand{\GL}{\mathrm{GL}}                 % Groupe linéaire général
\newcommand{\SL}{\mathrm{SL}}                 % Groupe spécial linéaire

% ========================================
% TOPOLOGIE ALGÉBRIQUE
% ========================================

% --- Espaces classiques ---
\newcommand{\Sph}[1]{S^{#1}}                  % Sphère S^n
\newcommand{\Boule}[1]{D^{#1}}                % Boule D^n
\newcommand{\Tore}{T^2}                        % Tore
\newcommand{\RP}[1]{\mathbb{RP}^{#1}}         % Espace projectif réel
\newcommand{\CP}[1]{\mathbb{CP}^{#1}}         % Espace projectif complexe
\newcommand{\simplex}{\Delta}                  % Simplexe standard

% --- Groupes d'homotopie ---
\newcommand{\pifond}[2]{\pi_1(#1,#2)}         % π₁(X,x₀)
\newcommand{\pihomo}[2]{\pi_{#1} (#2)}         % πₙ(X)

% --- Homologie et cohomologie ---
\newcommand{\Hsing}[1]{H_{#1}}                % Homologie H_n
\newcommand{\Cohsing}[1]{H^{#1}}              % Cohomologie H^n
\newcommand{\Htilde}[1]{\widetilde{H}_{#1}}   % Homologie réduite

% --- Complexes de chaînes ---
\newcommand{\Csing}[1]{C_{#1}}                % Chaînes singulières C_n
\newcommand{\Zsing}[1]{Z_{#1}}                % Cycles Z_n
\newcommand{\Bsing}[1]{B_{#1}}                % Bords B_n

% --- Relations ---
\newcommand{\iso}{\cong}                       % Isomorphisme
\newcommand{\rel}{\,\text{rel}\,}             % Relatif à

% ========================================
% NOTATIONS DIVERSES
% ========================================
\newcommand{\xn}{(x_n)_{n\inI}}              % Suite indexée par I
\newcommand{\yn}{(y_n)_{n\inJ}}              % Suite indexée par J